\chapter{Considerações Parciais}
\label{cap:consideracoes_parciais}

% -----------------------------------------------------------------
% Objetivo do capítulo:
% Fechar o documento de qualificação retomando o problema, descrevendo
% o estado atual da pesquisa, indicando o cronograma de trabalhos até
% a defesa e elencando direções de trabalho futuro. A versão de defesa
% expande este capítulo em "Conclusão" com resultados, contribuições
% efetivamente entregues e limitações observadas.
% -----------------------------------------------------------------


\section{Síntese parcial da pesquisa}
\label{sec:sintese_parcial}

Esta dissertação investiga, no nível da qualificação, como combinar sinais operacionais de máquina e evidências estatísticas históricas de rompimentos para inferir a provável origem de eventos de rompimento em uma fábrica de fios esmaltados de cobre e alumínio. A pergunta foi enunciada na Seção~\ref{sec:problema_pesquisa} e desdobrada nos objetivos específicos da Subseção~\ref{subsec:objetivos_especificos}.

Até este ponto, foram construídos os elementos fundacionais da proposta. A fundamentação do domínio industrial, apresentada no Capítulo~\ref{cap:fundamentacao}, situa o problema do rompimento de fio nas tradições de qualidade e confiabilidade da manufatura contínua e descreve os dados disponíveis para análise. A revisão sistemática da literatura, conduzida segundo o protocolo \sigla{PBS} detalhado no Capítulo~\ref{cap:revisao_sistematica}, selecionou 11 trabalhos correlatos e identificou a lacuna que a proposta preenche, a saber, a integração de detecção de anomalias, análise estatística histórica segmentada por lote e fornecedor, fusão de evidências e explicabilidade em um único sistema aplicado a fabricação de fios esmaltados.

A arquitetura do WireSense foi formalizada no Capítulo~\ref{cap:proposta}, com a definição dos dois pilares independentes, das transformações que constituem o pipeline de dados e da estrutura de saída de cada pilar. O Pilar de Sinais Operacionais foi desenhado para operar por máquina, com \emph{Isolation Forest} treinado sobre o histórico de intervalos sem incidente. O Pilar de Histórico Estatístico recebeu sua formalização estatística, com a hora-máquina como métrica de exposição, a Taxa Relativa por Material como indicador padronizado por categoria e o cruzamento de falhas por máquina como sinal complementar.

Os elementos ainda em construção até a defesa, descritos na Seção~\ref{sec:cronograma}, completam a proposta com a regra de integração dos dois pilares, a avaliação sobre dados históricos da fábrica e a discussão das contribuições efetivamente entregues. Esta qualificação não traz resultados experimentais, que constituem objeto do trabalho previsto entre esta etapa e a defesa.


\section{Cronograma para a defesa}
\label{sec:cronograma}

O cronograma de trabalhos previstos entre a qualificação e a defesa está organizado em seis etapas sequenciais, apresentadas na Tabela~\ref{tab:cronograma}. As etapas combinam atividades de pesquisa, implementação, experimentação e redação, e cobrem o período subsequente à qualificação.

\begin{table}[htbp]
    \centering
    \footnotesize
    \caption{Cronograma de trabalhos previstos até a defesa.}
    \label{tab:cronograma}
    \begin{tabularx}{\textwidth}{Xl}
        \toprule
        \textbf{Etapa}                                                                 & \textbf{Prazo previsto} \\
        \midrule
        Conclusão da revisão sistemática (extração e síntese dos trabalhos correlatos) & Ago--Set 2026 \\
        Refinamento da implementação dos dois pilares e da regra de integração          & Out--Nov 2026 \\
        Condução dos experimentos sobre dados históricos da fábrica                     & Dez 2026 -- Fev 2027 \\
        Validação com casos documentados de devoluções e de paradas internas            & Mar--Abr 2027 \\
        Redação do capítulo de Resultados e Discussão                                   & Mai--Jun 2027 \\
        Revisão geral do texto e preparação da defesa                                   & Jul 2027 \\
        \bottomrule
    \end{tabularx}
    \fonte{Elaborado pelo autor (\the\year).}
\end{table}

A primeira etapa concentra-se em encerrar a revisão sistemática com o aprofundamento das discussões narrativas e da síntese das lacunas identificadas. A segunda etapa avança a implementação dos pilares para além do estado descrito no Capítulo~\ref{cap:proposta}, incorporando ajustes de engenharia e a finalização da regra de integração. A terceira e a quarta etapa correspondem ao trabalho empírico sobre os dados históricos da fábrica, com avaliação sistemática do desempenho dos pilares e da regra de integração contra eventos documentados. As duas últimas etapas tratam da redação final, incluindo a expansão deste capítulo em uma versão de defesa que substitui a presente síntese parcial por um capítulo de Resultados e Discussão e por uma conclusão consolidada.

Para a quarta etapa, a fábrica dispõe de pelo menos um lote com causa documentada por inspeção pericial. Trata-se de um lote de cobre identificado pela equipe de qualidade como responsável por rompimentos sistemáticos durante o consumo, atribuídos à presença de inclusões metálicas no vergalhão e expostas durante a trefilação. Esse caso constitui o ponto de partida confirmatório da validação: a expectativa é que o WireSense classifique esse lote como evidência forte de matéria-prima, com taxa relativa ao lote acima do esperado e cruzamento de falhas distribuído entre as máquinas que o processaram. A partir desse caso, a validação se estende para os demais eventos documentados de devoluções de cliente e de paradas internas, com a expectativa de cobrir as classes de resposta previstas pela regra de integração apresentada na Seção~\ref{sec:regras_integracao}.


\section{Trabalhos futuros}
\label{sec:trabalhos_futuros}

Para além da defesa desta dissertação, identificam-se direções de extensão da pesquisa que ampliam o alcance do método proposto e o aproximam da operação cotidiana da fábrica.

A direção principal projeta a construção de uma ferramenta visual de apoio ao operador e à equipe técnica, que tomará o método proposto como núcleo de inferência e adicionará camadas de operação em tempo real, alertas acionáveis e integração com os sistemas de chão de fábrica. O gatilho previsto para a inferência em tempo real é o próprio registro operacional do evento. Assim que o operador anota no sistema da máquina que a parada foi causada por rompimento, o WireSense é acionado e produz a classificação da provável origem com base na janela sensorial imediatamente anterior ao evento e no histórico estatístico do lote em consumo, devolvendo o diagnóstico à interface operacional antes da retomada da produção. Essa integração transforma a análise hoje retrospectiva, indicada na Seção~\ref{sec:caracterizacao_pesquisa}, em ferramenta de apoio à decisão durante o turno, oferecendo à equipe de manutenção indicação imediata sobre se a próxima ação deve recair sobre o lote em consumo ou sobre o equipamento. O detalhamento dessa ferramenta, incluindo escolhas de interface, tecnologia e fluxo operacional, constitui agenda de trabalho posterior à conclusão do mestrado.

A ampliação da cobertura de sensores e de tipos de defeito é uma segunda direção de extensão. O Pilar de Sinais Operacionais foi desenhado sobre um subconjunto de variáveis priorizadas, e o Pilar de Histórico Estatístico opera sobre os rompimentos. A incorporação de defeitos de bolha, asperidade e instabilidade de forno como categorias adicionais de evento e a inclusão de variáveis sensoriais hoje fora do conjunto selecionado permitiriam estender a inferência da provável origem para outros sintomas operacionais relevantes.

A incorporação de realimentação de especialistas como fonte de rótulo fraco é uma terceira direção. Quando a equipe técnica de qualidade revisa um evento e atribui causa provável, esse julgamento pode ser registrado e usado como rótulo de baixa frequência mas alta confiança, refinando ao longo do tempo a calibração dos pilares e da regra de integração.

A generalização da arquitetura para outros processos de manufatura contínua com falhas multifatoriais é a quarta direção. A separação em pilares independentes, com fusão por regra determinística, é em princípio agnóstica ao domínio. Aplicações em fundição, laminação, extrusão e em outros processos com características semelhantes ao da fábrica analisada poderiam validar a transferibilidade do método proposto, com adaptação dos critérios de seleção de variáveis sensoriais e da estrutura categórica de matérias-primas.
